\BeleskeID{07-s019}
% element: 07-s019:d01
% element: 07-s019:e01
Razvoj Veriloga vezuje se za Gateway Design Automation sredinom 1980-ih. Ubrzo po uvođenju postao je de fakto standard za ASIC simulacione biblioteke. Široka rasprostranjenost biblioteka dala mu je prednost u dostupnosti simulacionih modela; VHDL je u opisanom poređenju imao manju podršku za širok raspon ASIC uređaja. Oba jezika koriste se i za simulaciju i za sintezu.

% element: 07-s019:e02
VITAL uvodi u VHDL standardni metod zapisa kašnjenja sličan onome u Verilogu. Time dobavljačima ASIC-a i simulacionih modela olakšava izradu VHDL modela iz postojećih Verilog biblioteka.

% element: 07-s019:e03
PLI je interfejs programskog jezika definisan u standardu Open Verilog International. Omogućava da se deo modela izvede izvan Veriloga, na primer u C-u, kada je tako model brži ili kada bi određena funkcija bila teško ili neefikasno izražena direktno u HDL-u.

% element: 07-s019:e04
VHDL uz IEEE standard 1076 nudi upravljanje dizajnom višeg nivoa: deklaracije konfiguracije i biblioteke. Te osobine odgovaraju velikim projektima. U većini razmatranih funkcija dva jezika imaju istu ulogu; različita sintaksa ne menja osnovne koncepte opisa hardvera.
